\section{反应与目标态布居}\label{sec:5-1}

在低能核结构与在束 $\gamma$ 谱学研究中，实验测量的终极目标是提取原子核特定量子态的自旋宇称 $I^\pi$、单粒子组态、激发能、电磁跃迁几率 $B(XL)$ 以及内禀形变参数 $(\beta, \gamma)$。然而，这些微观结构物理量并不能脱离核反应过程而直接孤立观测。探测器所记录的谱线强度、角分布各向异性与多普勒线形，在因果链条的源头均受到核反应激发机制与入射能区布居截面的严密制约。

从核结构问题出发选择核反应体系，构成了实验设计的逻辑起点。不同的核反应机制具有截然不同的角动量输入能力、能量沉积特性与量子相空间选择定则。若未对复合核形成、粒子蒸发退激、侧向馈入（Side-Feeding）以及竞争反应道截面建立定量且清晰的物理认识，容易在后续谱学分析中陷入两类方法学陷阱：其一，将未观测到弱跃迁或高自旋态武断解释为结构模型的失效，而忽视了反应布居窗口的本征截断；其二，在能级寿命与电磁跃迁几率反演中忽略未知滞后馈入的时间延迟效应，导致提取出的约化矩阵元产生系统性畸变。本节系统讨论核反应机制的选择反演原则，重点解析重离子熔合蒸发反应的运动学、激发能与角动量耗散机制，对比库仑激发、转移反应、光核荧光与弱衰变探针的多途径互补特性，给出谱学可见性的物理参数框架，并阐明全反应产额向特定目标态布居传递时的偏差机制与跨节接口。

\subsection{核结构问题与反应机制的选择反演}\label{subsec:5-1-1}

实验谱仪中测得的任何一条特征 $\gamma$ 跃迁，其被观测到的前提是原子核的母态量子态在反应中被有效布居。因此，核反应的设计是一个以待研究的微观结构自由度为目标的反演决策过程：
\begin{enumerate}
\item 高自旋集体带与转动阻尼研究：要求向原子核注入高角动量，同时将复合核的激发能控制在粒子蒸发残留的低激发区，使得核退激路径迅速收敛至冷核态附近的 Yrast 转动带或近 Yrast 轴对称/三轴摆动带。这通常选用重离子熔合蒸发反应（Heavy-Ion Fusion-Evaporation）。
\item 单粒子轨道与核壳演化研究：目标态通常具有准粒子激发或单粒子-核孔组态特征。此时强集体熔合反应的多步蒸发级联反而会掩盖特定的轨道跃迁，应当选用直接转移反应（Transfer Reactions），利用动量匹配选择定则选择性布居特定轨道角动量 $L$ 的单核子态，并通过转移截面角分布提取光谱因子（Spectroscopic Factor）。
\item 偶偶核低自旋集体振动与形状共存态：若需要测量低激发 $2^+_1, 4^+_1$ 或第 II 类极小 $0^+_2$ 态的电磁跃迁几率与实验室谱学四极矩，基于电磁相互作用的低能或多步库仑激发（Coulomb Excitation）能够大幅抑制强相互作用核力贡献，是提取集体电四极强度的重要探针。
\item 低自旋电磁模式与核对称能敏感态：针对剪刀振动模式（Scissors Mode）、微矮偶极共振（PDR）等低自旋态的布居，光核共振荧光（NRF）提供了从基态 $0^+$ 直接吸收单光子跃迁的高选择性窗口；同时放射性衰变亦可在低自旋区间提供互补约束。
\end{enumerate}
上述对比表明，不存在一种能够普遍通适于全部核结构问题的单一反应途径。每一种反应机制都在特定的自旋、同位旋、激发能与选择定则窗口内工作，并引入各自的观测偏置。

\subsection{重离子熔合蒸发反应与角动量输入}\label{subsec:5-1-2}

重离子熔合蒸发反应是在束高自旋态谱学研究的主力实验途径。在目标靶静止、非相对论两体完全融合运动学模型中，考虑质量为 $M_p$、实验室系动能为 $E_{\mathrm{lab}}$ 的入射束流离子轰击质量为 $M_t$ 的靶核，质心系有效碰撞能量为：
\begin{equation}\label{eq:5-1-ecm}
E_{\mathrm{cm}} = E_{\mathrm{lab}} \frac{M_t}{M_p + M_t}.
\end{equation}
当两核发生完全融合形成复合核（Compound Nucleus, CN）时，反应质量亏损 $Q = (M_p + M_t - M_{\mathrm{CN}})c^2$ 转化为内部激发，复合核的初始激发能满足：
\begin{equation}\label{eq:5-1-estar}
E^* = E_{\mathrm{cm}} + Q.
\end{equation}
入射离子在穿透固体靶晶格时产生阻止本领能损，在有限靶厚内引起束流能量沉积展宽（$\Delta E_{\mathrm{target}} = \int_0^x [-\mathrm{d}E/\mathrm{d}x']\mathrm{d}x'$），使得反应截面在靶内呈现连续能量积分效应。

在单势垒穿透模型中，熔合反应由各个分波的穿透概率决定。采用倒置抛物线势垒近似，Wong 导出了平滑单势垒熔合截面公式~\cite{book-bertulani-nuclear-reactions, book-lin-heavy-ion-reactions}：
\begin{equation}\label{eq:5-1-wong}
\sigma_{\mathrm{cap}}^{\mathrm{Wong}}(E_{\mathrm{cm}}) = \frac{\hbar\omega_B R_B^2}{2E_{\mathrm{cm}}}\ln\left[1+\exp\left(\frac{2\pi(E_{\mathrm{cm}}-V_B)}{\hbar\omega_B}\right)\right],
\end{equation}
这里 $R_B$ 是所取 $s$ 波总相互作用势垒的峰位置，$V_B$ 为该处势垒高度，$\hbar\omega_B$ 表征局部倒抛物线曲率。势垒位置不应无条件等同几何接触半径。

当能量高于势垒（$E_{\mathrm{cm}} > V_B$）且高分波近似满足锐截止（Sharp-Cutoff）条件时，参与熔合的最大无量纲轨道角动量分波 $\ell_{\mathrm{crit}}$ 满足：
\begin{equation}\label{eq:5-1-lcrit}
\ell_{\mathrm{crit}}^2 \approx \frac{2\mu R_B^2}{\hbar^2}(E_{\mathrm{cm}} - V_B),
\end{equation}
式中 $\mu = M_p M_t / (M_p + M_t)$ 为约化质量，轨道角动量大小为 $L \approx \hbar \ell$。复合核的总角动量 $\mathbf{J}$ 由入射轨道角动量与两核本征自旋矢量耦合而成：
\begin{equation}\label{eq:5-1-j-coupling}
\mathbf J_{\mathrm{CN}}=\mathbf L+\mathbf I_p+\mathbf I_t.
\end{equation}
上式各矢量均有角动量量纲，$\ell$ 是相应无量纲轨道量子数。Wong式在指定单势垒与连续分波近似下给出俘获型结果；将俘获视为完全熔合还需相应吸收或复合核形成假设。一般地，复合核形成包含 $P_{\mathrm{CN}}$，蒸发剩余产物还包含后续 $W_{\mathrm{sur}}$。多通道耦合、形成与存活的分工见第~4.8~节。

对零自旋弹靶和选定束轴，经典轨道角动量 $\mathbf L=\mathbf r\times\mathbf p$ 垂直于相对入射动量，无量纲量为 $\boldsymbol\ell=\mathbf L/\hbar$。相应初始量子描述可取 $m_\ell=0$\cite{diamond-1966-nuclear-alignment-heavy-ion-reactions}。非零入射自旋、粒子蒸发和后续级联会改变目标态的排列；实际角分布分析仍须标定取向与去取向，不能从入口几何固定所有态的统计张量符号。

在复合核形成后，系统通过轻粒子蒸发进行退激。每个蒸发出的轻粒子（$n, p, \alpha$）带走与其分离能和动能相当的能量（约 $S_n + 2T$），但通常仅带走有限的角动量，使残核在最后一个粒子蒸发阈能之下依然保留可观的自旋。随后，残核经由高能级密度连续谱区的统计 $\gamma$ 级联冷却至 Yrast 线附近，并注入离散能级序列。由此可见，入射轨道分波 $\ell$、复合核自旋 $J$ 与最终观测到的残核激发能级自旋 $I$ 属于反应演化不同阶段的物理量，不可混淆。

\subsection{多途径互补：库仑激发、转移反应、光核荧光与衰变}\label{subsec:5-1-3}

虽然重离子熔合蒸发是布居高自旋态的重要手段，但在完整构建原子核微观结构图景时，多类辅助探针提供了不可替代的互补观测窗口：
\begin{enumerate}
\item 库仑激发（Coulomb Excitation）：
当入射离子的碰撞轨迹保持在核相互作用力程之外、满足经典最近距离要求时，核力相互作用贡献受到强烈抑制，原子核主要受到时变电磁场的微扰激发~\cite{book-bertulani-nuclear-reactions}。需要澄清，势垒以下散射并不能保证核力尾部在数学上严格为零，实际实验通常通过散射角与能量的选择将核干涉效应控制在误差允许范围之内。库仑激发振幅可按电磁微扰阶数展开：
\begin{equation}\label{eq:5-1-coulex-amp}
a_{i \to f} = a_{i \to f}^{(1)} + a_{i \to f}^{(2)} + \dots, \qquad P_{i \to f} = |a_{i \to f}|^2 \approx |a_{i \to f}^{(1)}|^2 + 2\mathrm{Re}\left(a_{i \to f}^{(1)*} a_{i \to f}^{(2)}\right) + \dots
\end{equation}
一阶项给出跃迁几率正比于初末态之间的约化跃迁几率 $B(E\lambda; i \to f)$。在多步库仑激发中，二阶微扰振幅包含途径 $i \to f \to f$ 的重取向效应（Reorientation Effect），该项耦合了末态的对角矩阵元 $\langle I_f \| \mathcal{M}(E2) \| I_f \rangle$。必须严格区分物理概念：该对角矩阵元直接对应的是实验室系谱学静电四极矩 $Q_s(I_f)$，而非内禀坐标系的几何四极矩 $Q_0$。由 $Q_s$ 推断内禀电荷不对称角 $\gamma$ 必须依赖具体的投影转子模型或微观波函数假设。此外，实验多极矩阵元的反演通常依赖多步耦合通道程序（如 GOSIA）对多条跃迁路径与已知分支比进行联合拟合，并非单次无模型的直接测量。
\item 核子转移反应（Transfer Reactions）：
如 $(d, p)$、$(p, d)$、$(t, \alpha)$ 或重离子单核子剥离反应，其物理机制是直接在入射粒子与靶核之间转移单个或对核子。转移截面依赖于转移轨道角动量 $L$ 的动量匹配条件，其微分截面正比于单粒子光谱因子 $C^2 S$。转移反应选择性布居单粒子或单空穴成分纯净的能级，是检验壳模型平均场单粒子轨道占据与反转岛侵入态的关键工具。
\item 光核共振荧光（NRF）：
利用实光子束照射靶核，光子吸收遵循纯电磁跃迁选择定则。从偶偶核基态 $0^+$ 出发，实光子偶极吸收直接激发 $1^-$（$E1$）或 $1^+$（$M1$）态；电四极吸收亦允许直接布居 $2^+$（$E2$）态。NRF 反应在零轨道角动量转移条件下，对核介质中的磁偶极剪刀模式振动（Scissors Mode）与微矮电偶极共振具有极高的灵敏度。
\item 放射性 $\beta$ 衰变：
母核通过弱相互作用退激至子核，受到 $Q_\beta$ 衰变能窗以及 Fermi 与 Gamow-Teller 选择定则的制约。需要指出，$\beta$ 衰变母核并不普遍处于低自旋态；高自旋同核异能态母核（例如高自旋同质异能素）完全能够通过允许跃迁或禁戒跃迁布居子核中较高自旋的能级。布居态的自旋范围由母核实际自旋宇称与衰变算符性质决定，随后通过离散 $\gamma$ 级联向下退激。
\end{enumerate}
在科学论证中，不能将重离子熔合蒸发简单叙述为布居核激发态的排他路径，而应根据待测目标态的构型特征，结合多类探针在各自观测窗口内的互补优势。

\subsection{截面、阈能、丰度与谱学可见性边界}\label{subsec:5-1-4}

一条理论预言的 $\gamma$ 跃迁能否在在束探测阵列谱图上被辨识并达到具有统计显著性的信噪比，受到核反应截面、同位素丰度与探测仪器效率的联合制约。

在反应能量设定上，必须明确区分反应运动学阈能与库仑势垒：
\begin{itemize}
\item 运动学阈能（Kinematic Threshold）：对于吸热核反应（$Q < 0$），由四维动量守恒导出的发生该反应所需的最低实验室系动能为 $E_{\mathrm{th}} = -Q (M_p + M_t) / M_t$；放热反应（$Q > 0$）的运动学阈能为零；
\item 库仑势垒高度（Coulomb Barrier）：由两核相互接触时的库仑排斥势 $V_B$ 决定。对于带电粒子反应，即便体系能量高于运动学阈能，若质心能量远低于势垒，截面亦会因库仑穿透几率而发生指数级衰减。
\end{itemize}

设束流粒子数流率为 $I_{\mathrm{beam}}$（$\mathrm{s}^{-1}$），目标核素的面数密度为 $n_t$（$\mathrm{cm}^{-2}$），所用布居截面 $\sigma_{\mathrm{pop}}(i)$ 包含声明的直接及馈入途径。在薄靶、束流能量与响应可按所用能区平均的近似下，若 $b_{\gamma,k}=\lambda_{\gamma,k}/\lambda_{\mathrm{tot}}$ 是每次该态退激发射目标光子的概率，而 $\mathcal A_{\gamma,k}$ 是发射该光子条件下探测、保存、符合门与清洗的联合接受概率，则峰计数率为：
\begin{equation}\label{eq:5-1-yield-gamma}
Y_{\gamma,k}=I_{\mathrm{beam}}\,n_t\,\sigma_{\mathrm{pop}}(i)\,b_{\gamma,k}\,\mathcal A_{\gamma,k},
\end{equation}
若由面质量密度、丰度和摩尔质量求 $n_t$，应先明确这些量对应整个靶材、目标元素还是已分离的同位素，避免重复乘丰度。$\mathcal A_{\gamma,k}$ 可用经核验的响应模型分解；单光子效率、联合live与clean接受度若有相关性，不能直接作为独立边缘概率相乘。若求总计数，还需对束流与能量随时间的变化积分。

该公式确立了决定在束谱学“可见性（Visibility）”的核心物理参数空间：
\begin{enumerate}
\item 激发函数形态与束能优化：部分反应截面 $\sigma_{\mathrm{level}}(E)$ 强烈依赖于束流能量。蒸发道截面在近垒或过垒能区开启后上升并达到极大，随后由于竞争蒸发道（如发射更多核子或带电粒子道）开启而下降。激发函数的峰值半高宽并非普遍为几兆电子伏，其展宽尺度取决于反应两体运动学与蒸发相空间。实验中选取入射束能并不必须追求单道截面的数学极大值，而是必须综合权衡杂质反应道污染、目标自旋窗口分布、伴随强背景发射以及靶材热损伤极限。
\item 同位素丰度与杂质背景抑制：若采用天然丰度较低的单质靶（$f_{\mathrm{iso}} \ll 1$），伴随的非目标同位素将产生强烈的伴生反应通道，其发射的特征 $\gamma$ 射线与连续康普顿散射会严重压低目标弱信号的信噪比。
\item 阵列全能峰探测效率：单晶高纯锗探测器的光电吸收效率随能量升高而迅速衰减。在兆电子伏能区，即使反应产额可观，较低的探测效率亦会显著延长积累统计所需的时间。
\end{enumerate}

\subsection{全反应产额、侧向馈入与反演偏差传递}\label{subsec:5-1-5}

在利用测量谱线反演微观电磁矩阵元与能级寿命时，必须建立严密的统计区分：全反应道产额不等于目标能级的瞬时布居截面。

在重离子熔合蒸发中，一个特定的低自旋或中自旋能级 $|i\rangle$（如带内 $4^+$ 态），其总激发通量实际上是由两个时间尺度截然不同的物理通路共同构成的：
\begin{equation}\label{eq:5-1-feeding-split}
\Phi_{\mathrm{tot}}(i) = \Phi_{\mathrm{direct}}(i) + \Phi_{\mathrm{discrete}}(i) + \Phi_{\mathrm{side}}(i),
\end{equation}
其中：
\begin{itemize}
\item $\Phi_{\mathrm{direct}}(i)$：从粒子蒸发终点残核激发态经快速统计级联直接注入能级 $|i\rangle$ 的瞬发直接布居份额；
\item $\Phi_{\mathrm{discrete}}(i)$：沿本能级上方已知的离散能级序列逐级退激的离散级联馈入（Cascade Feeding）；
\item $\Phi_{\mathrm{side}}(i)$：从高能准连续谱区、非共振侧带或尚未分辨的微弱侧向结构注入该态的侧向馈入（Side-Feeding）。
\end{itemize}

取归一化注入时间概率密度 $f_{\mathrm{feed}}(t)$。在单一、常数总寿命 $\tau$ 的孤立态模型中，以下 $S(t)$ 是每次注入后离开该态的退激时间概率密度；布居概率为 $n_i(t)=\tau S(t)$，不将二者混称：
\begin{equation}\label{eq:5-1-feeding-conv}
S(t) = \int_0^t f_{\mathrm{feed}}(t') \frac{e^{-(t - t')/\tau}}{\tau} \, \mathrm{d}t'.
\end{equation}
忽略侧向馈入的时间延迟将对反演产生系统性偏置。在在束实验中，若侧向结构经过较慢的连续谱多极跃迁或微弱同核异能态注入，测得的时间衰变线形（如 DSAM 中的多普勒展宽多峰或反冲距离法中的未移位/移位峰比例）将呈现滞后尾部。若分析中武断假定为瞬时直接布居，测得的表观衰变时间将被延长，通常导致反演得到的平均寿命 $\tau$ 产生正向偏差（系统性偏大）；根据第~5.9~节中建立的反演公式 $B(E2) \propto 1/(\tau E_\gamma^5)$，高估的能级寿命将导致微观约化跃迁几率 $B(E2)$ 被系统性低估。但需要注意，偏置的具体方向与幅度取决于实际馈入时间谱形状与寿命测量时间窗口的相对重叠。

在待测态上方的离散跃迁设置门窗，可以选择某些已知馈入路径、减少部分未解析馈入的不确定性；它不自动消除该路径上各态的等待时间或全部慢馈入。门选和馈入模型应共同进入第~5.8~节的测量响应。
\begin{itemize}
\item 第 5.1 节确立的核反应输入与馈入网络，决定了第~5.2~节探测器所接收的原始波形流与第~5.6~节符合矩阵中的非对称本底结构；
\item 在第~5.8~节展开的在束寿命测量中，必须通过符合门选技术或参数化馈入时间函数模型显式处理侧向馈入响应；
\item 第 5.1 节阐明的反应选择性与通道可见性，为宏观谱学计数与微观结构动力学（第 4 章）的比对提供了物理边界。
\end{itemize}
